% 文献元数据核验日期：2026-09-10。
% DOI 注册元数据与出版社页面的核验记录见 ../文献与模型对应.md。
% 本文件直接使用 % 文献元数据核验日期：2026-09-10。
% DOI 注册元数据与出版社页面的核验记录见 ../文献与模型对应.md。
% 本文件直接使用 % 文献元数据核验日期：2026-09-10。
% DOI 注册元数据与出版社页面的核验记录见 ../文献与模型对应.md。
% 本文件直接使用 % 文献元数据核验日期：2026-09-10。
% DOI 注册元数据与出版社页面的核验记录见 ../文献与模型对应.md。
% 本文件直接使用 \input{sections/references}，无需 BibTeX。
\begingroup
\small\linespread{1.15}\selectfont
\begin{thebibliography}{99}
\setlength{\itemsep}{4pt}
\setlength{\parsep}{0pt}

\bibitem{kong2016}
Kong X, Bai L, Hu Q, et al.
Day-ahead optimal scheduling method for grid-connected microgrid based on energy storage control strategy[J].
Journal of Modern Power Systems and Clean Energy, 2016, 4(4): 648--658.
DOI: \url{https://doi.org/10.1007/s40565-016-0245-0}.

\bibitem{silva2020}
Silva V A, Aoki A R, Lambert-Torres G.
Optimal day-ahead scheduling of microgrids with battery energy storage system[J].
Energies, 2020, 13(19): 5188.
DOI: \url{https://doi.org/10.3390/en13195188}.

\bibitem{abunima2022}
Abunima H, Park W-H, Glick M B, et al.
Two-stage stochastic optimization for operating a renewable-based microgrid[J].
Applied Energy, 2022, 325: 119848.
DOI: \url{https://doi.org/10.1016/j.apenergy.2022.119848}.

\bibitem{marino2018}
Marino C, Quddus M A, Marufuzzaman M, et al.
A chance-constrained two-stage stochastic programming model for reliable microgrid operations under power demand uncertainty[J].
Sustainable Energy, Grids and Networks, 2018, 13: 66--77.
DOI: \url{https://doi.org/10.1016/j.segan.2017.12.007}.

\bibitem{grover2018}
Grover-Silva E, Heleno M, Mashayekh S, et al.
A stochastic optimal power flow for scheduling flexible resources in microgrids operation[J].
Applied Energy, 2018, 229: 201--208.
DOI: \url{https://doi.org/10.1016/j.apenergy.2018.07.114}.

\bibitem{palma2013}
Palma-Behnke R, Benavides C, Lanas F, et al.
A microgrid energy management system based on the rolling horizon strategy[J].
IEEE Transactions on Smart Grid, 2013, 4(2): 996--1006.
DOI: \url{https://doi.org/10.1109/TSG.2012.2231440}.

\bibitem{silvente2015}
Silvente J, Kopanos G M, Pistikopoulos E N, et al.
A rolling horizon optimization framework for the simultaneous energy supply and demand planning in microgrids[J].
Applied Energy, 2015, 155: 485--501.
DOI: \url{https://doi.org/10.1016/j.apenergy.2015.05.090}.

\bibitem{li2016}
Li Z, Zang C, Zeng P, et al.
Combined two-stage stochastic programming and receding horizon control strategy for microgrid energy management considering uncertainty[J].
Energies, 2016, 9(7): 499.
DOI: \url{https://doi.org/10.3390/en9070499}.

\bibitem{jani2022}
Jani A, Karimi H, Jadid S.
Two-layer stochastic day-ahead and real-time energy management of networked microgrids considering integration of renewable energy resources[J].
Applied Energy, 2022, 323: 119630.
DOI: \url{https://doi.org/10.1016/j.apenergy.2022.119630}.

\bibitem{li2020}
Li J, Tan C, Ren Z, et al.
A two-stage optimal scheduling model of microgrid based on chance-constrained programming in spot markets[J].
Processes, 2020, 8(1): 107.
DOI: \url{https://doi.org/10.3390/pr8010107}.

\bibitem{rockafellar2000}
Rockafellar R T, Uryasev S.
Optimization of conditional value-at-risk[J].
The Journal of Risk, 2000, 2(3): 21--41.
DOI: \url{https://doi.org/10.21314/JOR.2000.038}.

\bibitem{elmachtoub2017spo}
Elmachtoub A N, Grigas P.
Smart ``Predict, then Optimize''[EB/OL].
arXiv:1710.08005, 2017 (v5 updated 2020-11-19) [2026-09-10].
\url{https://arxiv.org/abs/1710.08005}.
DOI: \url{https://doi.org/10.48550/arXiv.1710.08005}.

\end{thebibliography}
\endgroup
，无需 BibTeX。
\begingroup
\small\linespread{1.15}\selectfont
\begin{thebibliography}{99}
\setlength{\itemsep}{4pt}
\setlength{\parsep}{0pt}

\bibitem{kong2016}
Kong X, Bai L, Hu Q, et al.
Day-ahead optimal scheduling method for grid-connected microgrid based on energy storage control strategy[J].
Journal of Modern Power Systems and Clean Energy, 2016, 4(4): 648--658.
DOI: \url{https://doi.org/10.1007/s40565-016-0245-0}.

\bibitem{silva2020}
Silva V A, Aoki A R, Lambert-Torres G.
Optimal day-ahead scheduling of microgrids with battery energy storage system[J].
Energies, 2020, 13(19): 5188.
DOI: \url{https://doi.org/10.3390/en13195188}.

\bibitem{abunima2022}
Abunima H, Park W-H, Glick M B, et al.
Two-stage stochastic optimization for operating a renewable-based microgrid[J].
Applied Energy, 2022, 325: 119848.
DOI: \url{https://doi.org/10.1016/j.apenergy.2022.119848}.

\bibitem{marino2018}
Marino C, Quddus M A, Marufuzzaman M, et al.
A chance-constrained two-stage stochastic programming model for reliable microgrid operations under power demand uncertainty[J].
Sustainable Energy, Grids and Networks, 2018, 13: 66--77.
DOI: \url{https://doi.org/10.1016/j.segan.2017.12.007}.

\bibitem{grover2018}
Grover-Silva E, Heleno M, Mashayekh S, et al.
A stochastic optimal power flow for scheduling flexible resources in microgrids operation[J].
Applied Energy, 2018, 229: 201--208.
DOI: \url{https://doi.org/10.1016/j.apenergy.2018.07.114}.

\bibitem{palma2013}
Palma-Behnke R, Benavides C, Lanas F, et al.
A microgrid energy management system based on the rolling horizon strategy[J].
IEEE Transactions on Smart Grid, 2013, 4(2): 996--1006.
DOI: \url{https://doi.org/10.1109/TSG.2012.2231440}.

\bibitem{silvente2015}
Silvente J, Kopanos G M, Pistikopoulos E N, et al.
A rolling horizon optimization framework for the simultaneous energy supply and demand planning in microgrids[J].
Applied Energy, 2015, 155: 485--501.
DOI: \url{https://doi.org/10.1016/j.apenergy.2015.05.090}.

\bibitem{li2016}
Li Z, Zang C, Zeng P, et al.
Combined two-stage stochastic programming and receding horizon control strategy for microgrid energy management considering uncertainty[J].
Energies, 2016, 9(7): 499.
DOI: \url{https://doi.org/10.3390/en9070499}.

\bibitem{jani2022}
Jani A, Karimi H, Jadid S.
Two-layer stochastic day-ahead and real-time energy management of networked microgrids considering integration of renewable energy resources[J].
Applied Energy, 2022, 323: 119630.
DOI: \url{https://doi.org/10.1016/j.apenergy.2022.119630}.

\bibitem{li2020}
Li J, Tan C, Ren Z, et al.
A two-stage optimal scheduling model of microgrid based on chance-constrained programming in spot markets[J].
Processes, 2020, 8(1): 107.
DOI: \url{https://doi.org/10.3390/pr8010107}.

\bibitem{rockafellar2000}
Rockafellar R T, Uryasev S.
Optimization of conditional value-at-risk[J].
The Journal of Risk, 2000, 2(3): 21--41.
DOI: \url{https://doi.org/10.21314/JOR.2000.038}.

\bibitem{elmachtoub2017spo}
Elmachtoub A N, Grigas P.
Smart ``Predict, then Optimize''[EB/OL].
arXiv:1710.08005, 2017 (v5 updated 2020-11-19) [2026-09-10].
\url{https://arxiv.org/abs/1710.08005}.
DOI: \url{https://doi.org/10.48550/arXiv.1710.08005}.

\end{thebibliography}
\endgroup
，无需 BibTeX。
\begingroup
\small\linespread{1.15}\selectfont
\begin{thebibliography}{99}
\setlength{\itemsep}{4pt}
\setlength{\parsep}{0pt}

\bibitem{kong2016}
Kong X, Bai L, Hu Q, et al.
Day-ahead optimal scheduling method for grid-connected microgrid based on energy storage control strategy[J].
Journal of Modern Power Systems and Clean Energy, 2016, 4(4): 648--658.
DOI: \url{https://doi.org/10.1007/s40565-016-0245-0}.

\bibitem{silva2020}
Silva V A, Aoki A R, Lambert-Torres G.
Optimal day-ahead scheduling of microgrids with battery energy storage system[J].
Energies, 2020, 13(19): 5188.
DOI: \url{https://doi.org/10.3390/en13195188}.

\bibitem{abunima2022}
Abunima H, Park W-H, Glick M B, et al.
Two-stage stochastic optimization for operating a renewable-based microgrid[J].
Applied Energy, 2022, 325: 119848.
DOI: \url{https://doi.org/10.1016/j.apenergy.2022.119848}.

\bibitem{marino2018}
Marino C, Quddus M A, Marufuzzaman M, et al.
A chance-constrained two-stage stochastic programming model for reliable microgrid operations under power demand uncertainty[J].
Sustainable Energy, Grids and Networks, 2018, 13: 66--77.
DOI: \url{https://doi.org/10.1016/j.segan.2017.12.007}.

\bibitem{grover2018}
Grover-Silva E, Heleno M, Mashayekh S, et al.
A stochastic optimal power flow for scheduling flexible resources in microgrids operation[J].
Applied Energy, 2018, 229: 201--208.
DOI: \url{https://doi.org/10.1016/j.apenergy.2018.07.114}.

\bibitem{palma2013}
Palma-Behnke R, Benavides C, Lanas F, et al.
A microgrid energy management system based on the rolling horizon strategy[J].
IEEE Transactions on Smart Grid, 2013, 4(2): 996--1006.
DOI: \url{https://doi.org/10.1109/TSG.2012.2231440}.

\bibitem{silvente2015}
Silvente J, Kopanos G M, Pistikopoulos E N, et al.
A rolling horizon optimization framework for the simultaneous energy supply and demand planning in microgrids[J].
Applied Energy, 2015, 155: 485--501.
DOI: \url{https://doi.org/10.1016/j.apenergy.2015.05.090}.

\bibitem{li2016}
Li Z, Zang C, Zeng P, et al.
Combined two-stage stochastic programming and receding horizon control strategy for microgrid energy management considering uncertainty[J].
Energies, 2016, 9(7): 499.
DOI: \url{https://doi.org/10.3390/en9070499}.

\bibitem{jani2022}
Jani A, Karimi H, Jadid S.
Two-layer stochastic day-ahead and real-time energy management of networked microgrids considering integration of renewable energy resources[J].
Applied Energy, 2022, 323: 119630.
DOI: \url{https://doi.org/10.1016/j.apenergy.2022.119630}.

\bibitem{li2020}
Li J, Tan C, Ren Z, et al.
A two-stage optimal scheduling model of microgrid based on chance-constrained programming in spot markets[J].
Processes, 2020, 8(1): 107.
DOI: \url{https://doi.org/10.3390/pr8010107}.

\bibitem{rockafellar2000}
Rockafellar R T, Uryasev S.
Optimization of conditional value-at-risk[J].
The Journal of Risk, 2000, 2(3): 21--41.
DOI: \url{https://doi.org/10.21314/JOR.2000.038}.

\bibitem{elmachtoub2017spo}
Elmachtoub A N, Grigas P.
Smart ``Predict, then Optimize''[EB/OL].
arXiv:1710.08005, 2017 (v5 updated 2020-11-19) [2026-09-10].
\url{https://arxiv.org/abs/1710.08005}.
DOI: \url{https://doi.org/10.48550/arXiv.1710.08005}.

\end{thebibliography}
\endgroup
，无需 BibTeX。
\begingroup
\small\linespread{1.15}\selectfont
\begin{thebibliography}{99}
\setlength{\itemsep}{4pt}
\setlength{\parsep}{0pt}

\bibitem{kong2016}
Kong X, Bai L, Hu Q, et al.
Day-ahead optimal scheduling method for grid-connected microgrid based on energy storage control strategy[J].
Journal of Modern Power Systems and Clean Energy, 2016, 4(4): 648--658.
DOI: \url{https://doi.org/10.1007/s40565-016-0245-0}.

\bibitem{silva2020}
Silva V A, Aoki A R, Lambert-Torres G.
Optimal day-ahead scheduling of microgrids with battery energy storage system[J].
Energies, 2020, 13(19): 5188.
DOI: \url{https://doi.org/10.3390/en13195188}.

\bibitem{abunima2022}
Abunima H, Park W-H, Glick M B, et al.
Two-stage stochastic optimization for operating a renewable-based microgrid[J].
Applied Energy, 2022, 325: 119848.
DOI: \url{https://doi.org/10.1016/j.apenergy.2022.119848}.

\bibitem{marino2018}
Marino C, Quddus M A, Marufuzzaman M, et al.
A chance-constrained two-stage stochastic programming model for reliable microgrid operations under power demand uncertainty[J].
Sustainable Energy, Grids and Networks, 2018, 13: 66--77.
DOI: \url{https://doi.org/10.1016/j.segan.2017.12.007}.

\bibitem{grover2018}
Grover-Silva E, Heleno M, Mashayekh S, et al.
A stochastic optimal power flow for scheduling flexible resources in microgrids operation[J].
Applied Energy, 2018, 229: 201--208.
DOI: \url{https://doi.org/10.1016/j.apenergy.2018.07.114}.

\bibitem{palma2013}
Palma-Behnke R, Benavides C, Lanas F, et al.
A microgrid energy management system based on the rolling horizon strategy[J].
IEEE Transactions on Smart Grid, 2013, 4(2): 996--1006.
DOI: \url{https://doi.org/10.1109/TSG.2012.2231440}.

\bibitem{silvente2015}
Silvente J, Kopanos G M, Pistikopoulos E N, et al.
A rolling horizon optimization framework for the simultaneous energy supply and demand planning in microgrids[J].
Applied Energy, 2015, 155: 485--501.
DOI: \url{https://doi.org/10.1016/j.apenergy.2015.05.090}.

\bibitem{li2016}
Li Z, Zang C, Zeng P, et al.
Combined two-stage stochastic programming and receding horizon control strategy for microgrid energy management considering uncertainty[J].
Energies, 2016, 9(7): 499.
DOI: \url{https://doi.org/10.3390/en9070499}.

\bibitem{jani2022}
Jani A, Karimi H, Jadid S.
Two-layer stochastic day-ahead and real-time energy management of networked microgrids considering integration of renewable energy resources[J].
Applied Energy, 2022, 323: 119630.
DOI: \url{https://doi.org/10.1016/j.apenergy.2022.119630}.

\bibitem{li2020}
Li J, Tan C, Ren Z, et al.
A two-stage optimal scheduling model of microgrid based on chance-constrained programming in spot markets[J].
Processes, 2020, 8(1): 107.
DOI: \url{https://doi.org/10.3390/pr8010107}.

\bibitem{rockafellar2000}
Rockafellar R T, Uryasev S.
Optimization of conditional value-at-risk[J].
The Journal of Risk, 2000, 2(3): 21--41.
DOI: \url{https://doi.org/10.21314/JOR.2000.038}.

\bibitem{elmachtoub2017spo}
Elmachtoub A N, Grigas P.
Smart ``Predict, then Optimize''[EB/OL].
arXiv:1710.08005, 2017 (v5 updated 2020-11-19) [2026-09-10].
\url{https://arxiv.org/abs/1710.08005}.
DOI: \url{https://doi.org/10.48550/arXiv.1710.08005}.

\end{thebibliography}
\endgroup
